\documentclass[11pt,a4paper]{article}
\usepackage[utf8]{inputenc}
\usepackage[T2A]{fontenc}
\usepackage[russian]{babel}
\usepackage[margin=2.2cm]{geometry}
\usepackage{amsmath,amssymb}
\usepackage{enumitem}
\setlist{itemsep=3pt,topsep=4pt}
\pagestyle{plain}

\begin{document}

\begin{center}
{\Large\bfseries Основные выводы}\\[4pt]
Никита Брусникин, 19.11.2026
\end{center}

\section*{1. Ванна Кондо в симуляторе}
\begin{itemize}
\item Добавлен микроскопический диссипатор: спины, связанные обменом с газом
  свободных электронов (размерность газа $d=1,2,3$), во втором порядке по
  $g=\rho J$ в секулярной форме Блоха--Рэдфилда. Параметры среды --- $g$,
  $k_BT$ и $k_Fd$ --- заменяют феноменологическую скорость $\gamma(t)$.
\item Динамика проверена: лиувиллиан построенных операторов коллапса
  воспроизводит скорости переходов и $1/T_2$ формулы статьи с точностью
  $10^{-14}$; стационарное состояние --- гиббсово.
\end{itemize}

\section*{2. Статья I: Delgado, Hirjibehedin, Fern\'andez-Rossier, Surf.\ Sci.\ 630, 337 (2014)}
\begin{itemize}
\item Воспроизведены: скорость релаксации одного спина и закон Корринги
  (функция $G(\Delta)$ совпадает с численной квадратурой до
  $2.6\cdot10^{-13}$), общая формула сдвига уровней ур.~21 ($2.2\cdot10^{-16}$),
  частный случай спина $1/2$ ур.~22.
\item Расхождение: ур.~21 даёт коэффициент сдвига щели $3/4$ для любого спина
  ($S=1,\tfrac32,2,\tfrac52$), а в частных формулах напечатано $3/2$ (ур.~24)
  и $3/16$ (ур.~23). Для уменьшения щели на 10~мэВ это даёт $g=0.28$ против
  1.4 по ур.~23.
\item Префактор скорости $\pi/2$ в этой статье против $\pi/4$ в статье II ---
  разные соглашения; на отношения скоростей не влияет.
\end{itemize}

\section*{3. Статья II: Delgado, Fern\'andez-Rossier, PRB 95, 075413 (2017)}
\begin{itemize}
\item Воспроизведены формулы для димера (рис.~2, до $6.7\cdot10^{-16}$),
  коллективное усиление спиновой волны $\Gamma_{SW1}(N)=N\,\Gamma_{SW1}(1)$
  (рис.~3, до $8\cdot10^{-15}$), дефазировка изинговских цепочек (рис.~4,
  прямой расчёт против суммы по ядру до $5.7\cdot10^{-14}$).
\item Напечатанные рисунки оцифрованы и сравнены с формулами:
  \begin{itemize}
  \item рис.~2(b): кривые 2D и 3D равны половине ур.~31--32 (медиана
    отклонения 0.001 и 0.003; от самих уравнений 0.45 и 0.48);
  \item рис.~4 построен по ур.~16 в напечатанном виде (медиана 0.005--0.018
    против 0.08--0.50 для трансляционно-инвариантной формы). Ур.~16 зависит
    от выбора начала координат; приложение~B статьи и наш код --- нет.
  \end{itemize}
\item В приложении C ур.~C4 не соответствует рисунку, ур.~C6 не сходится к
  заявленному пределу; правильная сумма выписана.
\end{itemize}

\section*{4. Два кубита с ванной Кондо}
\begin{itemize}
\item XXZ ($J=0.8$, $g=0.1$, $k_BT=0.2$): медиана $T_2^E$ меняется в 4.6 раза
  при изменении расстояния между кубитами; для одномерного газа 474 при
  $k_Fd=0$ и 104 при $k_Fd=\pi/2$, период $\pi$. Ширина 95\%-го интервала
  6--11\% медианы (30 начальных состояний на точку).
\item TFIM: изменение лишь в 1.14 раза, при 30 состояниях на точку не
  разрешено.
\item Секулярное приближение при этих параметрах выполнено: отношение
  скорости к расстоянию между боровскими частотами $\le0.0044$ (XXZ) и
  $\le0.094$ (TFIM).
\item $T_2$ в вычислительном базисе для XXZ плохо определено: стационарное
  (почти синглетное) состояние сохраняет $l_1$-когерентность 1, выше порога
  у 79\% начальных состояний. Цель $T_2^E$ этого недостатка не имеет.
\item Итог для проекта: появляется сценарий, в котором время декогеренции
  задано скрытым параметром среды (расстоянием $k_Fd$), а не
  феноменологической $\gamma(t)$.
\end{itemize}

\section*{5. Замечания А.~Е.~Теретенкова}
\begin{itemize}
\item Модель наблюдения с конечным числом повторов $N$ готова: оценки средних
  из выборки исходов, дисперсия $(1-\langle\sigma\rangle^2)/N$ подтверждена
  на 2000 повторениях; когерентность и чистота двух кубитов доступны только
  при томографии (9 настроек), а не по шести локальным средним. Оценки чистоты
  и когерентности по восстановленной $\hat\rho$ смещены вверх.
\item Симулятор записей последовательных проективных измерений с коллапсом
  готов и проверен (эффект Зенона, затухание $e^{-\gamma t}$ на 4000 записях).
\item Тексты статей исправлены: ансамблевая постановка вместо сценария
  ``предупреждение во время вычисления'', переписан раздел о входах двух
  кубитов, добавлены стоимость набора данных и работы из письма, исправлены
  числа E5 во введении статьи~1.
\item Не завершено: обучение и проверка моделей на входах с конечным $N$,
  сравнение ``шесть локальных средних против томографии'', таблица ``точность
  против стоимости набора''; модель и цель на записях последовательных
  измерений.
\end{itemize}

\section*{6. Что остаётся}
\begin{itemize}
\item Расчёты по замечаниям 1--3 (предыдущий пункт).
\item Сдвиг уровней (Lamb shift) в динамике и сравнение с несекулярным
  уравнением Блоха--Рэдфилда.
\item Ответы научного руководителя на вопросы о сценарии статей и о
  направлении последовательных измерений.
\end{itemize}

\end{document}
